\section{Introduction}

The evolution of wireless communication networks toward the sixth generation (6G) is driven by the demand for ultra-high data rates and massive connectivity. To meet these requirements, extremely large-scale multiple-input multiple-output (XL-MIMO) operating in the sub-terahertz (sub-THz) and THz bands has emerged as a pivotal technology \cite{lu2024tutorial, amiri2018extremely}. By deploying hundreds or thousands of antenna elements, XL-MIMO provides substantial spatial degrees of freedom and array gains. However, the significant increase in the antenna array aperture, coupled with high operating frequencies, alters the electromagnetic propagation characteristics. Specifically, the conventional far-field planar wave assumption becomes invalid, as users are frequently located within the Fresnel zone of the base station (BS). Consequently, the communication paradigm shifts from far-field planar waves to near-field spherical wave propagation, which necessitates a re-evaluation of transceiver design \cite{zhi2024performance, lu2024hybrid}.

A consequence of near-field spherical wave propagation in XL-MIMO systems is the emergence of spatial non-stationarity \cite{iimori2022grant}. Unlike conventional massive MIMO systems where the signal energy is uniformly distributed across the entire array, the signal from a specific user in an XL-MIMO system is effectively captured by only a localized subset of antennas. This localized subset is formally defined as the visibility region (VR) \cite{tang2024joint, xu2024joint}. The presence of VRs implies that processing the received signal across the entire antenna array is computationally inefficient and potentially detrimental to system performance. Traditional linear detection algorithms, such as zero-forcing (ZF) and minimum mean square error (MMSE) detectors, require large-scale matrix inversions. When applied to the full XL-MIMO array, these methods incur substantial computational complexity. Furthermore, incorporating antennas outside the user's VR into the detection process introduces noise accumulation, which degrades the overall detection accuracy \cite{zhi2024performance}.

To overcome the computational bottlenecks of traditional linear detectors, deep learning (DL) has been investigated for MIMO detection \cite{he2020model}. Among various DL architectures, graph neural networks (GNNs) have demonstrated effectiveness by modeling the MIMO detection problem as a message-passing process on a bipartite graph \cite{scotti2020graph, he2023message}. Existing GNN-based detectors achieve near-optimal performance with reduced complexity compared to iterative algorithms \cite{he2023gnn, li2022edge}. However, a technical challenge remains: state-of-the-art GNN detectors predominantly rely on fully-connected graph topologies, where every user node is connected to every antenna node. In XL-MIMO scenarios, this fully-connected assumption exhibits limited scalability, yielding a computational complexity of $\mathcal{O}(NK)$ per layer, where $N$ and $K$ are the numbers of antennas and users, respectively. As $N$ grows to the thousands, fully-connected GNNs encounter dimensionality challenges and memory overflow (out-of-memory, OOM) issues. Moreover, in wideband THz XL-MIMO systems, the near-field beam splitting effect causes the VR to shift across different subcarriers \cite{cui2022near, gao2024deep}, rendering static, fully-connected graph structures suboptimal and less capable of capturing the frequency-dependent spatial non-stationarity.

To address these challenges, this paper proposes Graph-VR-Det, a graph-based near-field detection scheme for XL-MIMO systems via dynamic VR discovery. By exploiting spatial non-stationarity, the XL-MIMO detection is modeled as a sparse message-passing problem, reducing computational overhead while maintaining detection precision. The contributions of this paper are summarized as follows:
\begin{itemize}
    \item A dynamic, pilot-based VR discovery algorithm is proposed to extract sparse user-antenna bipartite graph topologies without requiring prior channel state information (CSI). By evaluating the pilot-energy distributions, the algorithm identifies the valid VR edges for each user.
    \item A Sparse-GNN detector is developed, which restricts evidence processing and message passing exclusively to the discovered valid edges. This sparse activation mechanism reduces the computational complexity from scaling linearly with the total number of antennas to scaling linearly with the VR size.
    \item A frequency-aware feature mechanism is introduced to mitigate the wideband near-field beam splitting effect at THz frequencies. By integrating normalized frequency offsets into the edge features and constructing a shared graph topology across subcarriers, the proposed GNN compensates for frequency-dependent spatial shifts.
    \item Numerical evaluations demonstrate that the proposed Sparse-GNN achieves up to a 3.3-fold inference speedup over fully-connected networks. At an optimal sparsity level, the sparse graph topology induces a regularization effect that filters out noise, allowing the Sparse-GNN to surpass the bit error rate (BER) performance of dense networks. The architecture also exhibits zero-shot generalization capabilities, scaling to 1024-antenna arrays without retraining. However, it is noted that in wideband scenarios, the proposed method exhibits an error floor at high signal-to-noise ratios compared to genie-aided linear detectors, highlighting a trade-off between computational efficiency and high-SNR precision.
\end{itemize}

The remainder of this paper is organized as follows. Section II introduces the system model, detailing the near-field wideband channel and the problem formulation. Section III presents the proposed Graph-VR-Det methodology, including the dynamic VR discovery and the frequency-aware Sparse-GNN architecture. Section IV provides numerical results and complexity analyses. Finally, Section V concludes the paper.